% Part: first-order-logic
% Chapter: sequent-calculus
% Section: proving-things-quant

\documentclass[../../../include/open-logic-section]{subfiles}

\begin{document}

\olfileid[id]{fol}{seq}{prq}

\olsection{\usetoken{P}{derivation} dengan Kuantor}

\begin{ex}
Berikan sebuah $\Log{LK}$-!!{derivation} untuk sekuen
$\lexists[x][\lnot !A(x)] \Sequent \lnot \lforall[x][!A(x)]$.

Ketika menangani kuantor, kita harus memastikan bahwa syarat variabel
eigen tidak dilanggar, dan kadang-kadang hal ini mengharuskan kita mengubah
urutan pelaksanaan inferensi tertentu. Secara umum, sebaiknya kita mencoba
menangani terlebih dahulu aturan-aturan yang tunduk pada syarat variabel
eigen (aturan-aturan tersebut akan berada lebih bawah dalam pembuktian
akhir). Sebaiknya kita juga melihat beberapa langkah ke depan dan mencoba
menebak bentuk sekuen awalnya. Dalam kasus kita, bentuknya harus menyerupai
$!A(a) \Sequent !A(a)$. Artinya, ketika ``membalik'' aturan kuantor, kita
harus memilih suku yang sama---yang akan kita sebut $a$---untuk aturan
$\lforall$ maupun $\lexists$. Jika kita memilih suku yang berbeda untuk
setiap aturan, kita akan berakhir dengan bentuk seperti $!A(a) \Sequent
!A(b)$, yang tentu saja tidak dapat diturunkan.

Seperti biasa, kita mulai dengan menulis
\begin{prooftree}
\AxiomC{}
\UnaryInf$\lexists[x][\lnot !A(x)] \fCenter \lnot \lforall[x][!A(x)]$
\end{prooftree}
Kita dapat menerapkan aturan \LeftR{\exists} atau aturan
\RightR{\lnot}. Karena aturan \LeftR{\exists} tunduk pada syarat variabel
eigen, sebaiknya aturan itu ditangani lebih awal, jadi kita terapkan aturan
tersebut terlebih dahulu.
\begin{prooftree}
\AxiomC{}
\UnaryInf$ \lnot !A(a) \fCenter \lnot \lforall[x][!A(x)]$
\RightLabel{\LeftR{\lexists}}
\UnaryInf$ \lexists[x][\lnot !A(x)] \fCenter \lnot \lforall[x][!A(x)]$
\end{prooftree}
Dengan menerapkan aturan \LeftR{\lnot} dan \RightR{\lnot} secara mundur,
kita memperoleh
\begin{prooftree}
\AxiomC{}
\UnaryInf$\lforall[x][!A(x)] \fCenter !A(a)$
\RightLabel{\LeftR{\lnot}}
\UnaryInf$\lnot !A(a), \lforall[x][!A(x)] \fCenter $
\RightLabel{\LeftR{\Exchange}}
\UnaryInf$\lforall[x][!A(x)], \lnot !A(a) \fCenter $
\RightLabel{\RightR{\lnot}}
\UnaryInf$ \lnot !A(a) \fCenter \lnot \lforall[x] !A(x)$
\RightLabel{\LeftR{\lexists}}
\UnaryInf$ \lexists[x] \lnot !A(x) \fCenter \lnot \lforall[x] !A(x)$
\end{prooftree}
Pada tahap ini, satu-satunya pilihan kita adalah menerapkan aturan
\LeftR{\forall}. Karena aturan ini tidak tunduk pada pembatasan variabel
eigen, kita dapat melanjutkan. Ingat bahwa kita ingin memperoleh sekuen
awal (berbentuk $!A(a) \Sequent !A(a)$), jadi kita harus memilih $a$
sebagai argumen bagi $!A$ ketika menerapkan aturan tersebut.
\begin{prooftree}
\Axiom$!A(a) \fCenter !A(a)$
\RightLabel{\LeftR{\lforall}}
\UnaryInf$\lforall[x][!A(x)] \fCenter !A(a)$
\RightLabel{\LeftR{\lnot}}
\UnaryInf$\lnot !A(a), \lforall[x][!A(x)] \fCenter $
\RightLabel{\LeftR{\Exchange}}
\UnaryInf$\lforall[x][!A(x)], \lnot !A(a) \fCenter $
\RightLabel{\RightR{\lnot}}
\UnaryInf$ \lnot !A(a) \fCenter \lnot \lforall[x][!A(x)]$
\RightLabel{\LeftR{\lexists}}
\UnaryInf$ \lexists[x][ \lnot !A(x)] \fCenter \lnot \lforall[x][!A(x)]$
\end{prooftree}
Terutama ketika menangani kuantor, pada tahap ini penting untuk memeriksa
kembali bahwa syarat variabel eigen tidak dilanggar. Karena satu-satunya
aturan yang kita terapkan dan tunduk pada syarat variabel eigen adalah
\LeftR{\exists}, sedangkan variabel eigen~$a$ tidak muncul dalam sekuen
bawahnya (yakni sekuen akhir), ini merupakan !!a{derivation} yang benar.
\end{ex}

\begin{prob}
Berikan !!{derivation}s untuk sekuen-sekuen berikut:
\begin{enumerate}
\item $\Sequent (\lforall[x][!A(x)] \land \lforall[y][!B(y)]) \lif
\lforall[z][(!A(z) \land !B(z))]$.
\item $\Sequent (\lexists[x][!A(x)] \lor \lexists[y][!B(y)]) \lif
\lexists[z][(!A(z) \lor !B(z))]$.
\item $\lforall[x][(!A(x) \lif !B)] \Sequent \lexists[y][!A(y)] \lif !B$.
\item $\lforall[x][\lnot !A(x)] \Sequent \lnot\lexists[x][!A(x)]$.
\item $\Sequent \lnot\lexists[x][!A(x)] \lif \lforall[x][\lnot !A(x)]$.
\item $\Sequent \lnot\lexists[x][\lforall[y][((!A(x,y) \lif \lnot
!A(y,y)) \land (\lnot !A(y,y) \lif !A(x,y)))]]$.
\end{enumerate}
\end{prob}

\begin{prob}
Berikan !!{derivation}s untuk sekuen-sekuen berikut:
\begin{enumerate}
\item $\Sequent \lnot\lforall[x][!A(x)] \lif \lexists[x][\lnot!A(x)]$.
\item $(\lforall[x][!A(x)] \lif !B) \Sequent \lexists[y][(!A(y) \lif !B)]$.
\item $\Sequent \lexists[x][(!A(x) \lif \lforall[y][!A(y)])]$.
\end{enumerate}
(Semua sekuen ini memerlukan aturan~$\RightR{\Contraction}$.)
\end{prob}

\end{document}
